\chapter{Ensayos y resultados} % Main chapter title

\label{Chapter4}

En este capítulo se detalla, en primer lugar, cómo se evalúan los modelos estudiados en este trabajo. Luego, se explica cómo se evalúa el entorno de Airflow, con énfasis en los dos DAGs implementados. Posteriormente, se describen las pruebas realizadas en el servicio web y se presentan cada uno de sus \textit{endpoints}. Finalmente, se detalla una prueba realizada con otro dataset.

\definecolor{mygreen}{rgb}{0,0.6,0}
\definecolor{mygray}{rgb}{0.5,0.5,0.5}
\definecolor{mymauve}{rgb}{0.58,0,0.82}

%----------------------------------------------------------------------------------------
%	SECTION 1
%----------------------------------------------------------------------------------------
\section{Evaluación de modelos}
\label{sec:pruebas_indiv_modelos}

Una vez entrenados, se evalúan los modelos y sus versiones con el fin de encontrar al mejor de todos ellos. A continuación se detallan las métricas utilizadas como así también el proceso de comparación entre modelos.

\subsection{Métricas de modelos}
En este trabajo se estudian los modelos de Logistic Regression, SVM y XGBoost. Tal como se explica en la sección \ref{sec:flujo_trabajo}, en cada una de las iteraciones de estudio se genera una nueva versión de estos modelos. Esto aplica tanto a los modelos individuales como en los basados en \textit{pipelines}, tal como se describe en la sección \ref{sec:pipelines_y_modelos}.

Debido a la gran cantidad de versiones de modelos creados, surge la necesidad de definir métricas que permitan determinar cuál de estos es el que mejor permite predecir si una empresa entra en quiebra. En la tabla \ref{tab:metricas-de-modelos} se detallan las métricas utilizadas en este trabajo.

\begin{table}[H]
    \centering
    \caption{Métricas de modelos.}
    \begin{tabularx}{\textwidth}{l X X}
        \toprule
        \textbf{Métrica} & \textbf{¿Qué mide?} & \textbf{¿Por qué se eligió?} \\
        \midrule
        Recall & Capacidad para identificar todos los casos de quiebra reales. & Minimiza los falsos negativos en riesgos financieros. \\
        Precision & Exactitud de predecir quiebra. & Evita dar falsas alarmas. \\
        F1-Score & Balance entre Precision y Recall. & Evaluación general en una sola métrica. \\
        F2-Score & Balance que prioriza el Recall. & Prioriza detección de quiebras. \\
        AUPRC & Área bajo la curva de Precision y Recall. & Métrica robusta para datasets con clases minoritarias. \\
        ROC-AUC & Capacidad de discriminar entre clases. & Evalúa el modelo globalmente. \\
        \bottomrule
		\hline
    \end{tabularx}
    \label{tab:metricas-de-modelos}
\end{table}

En la tabla \ref{tab:resultados-metricas} se muestran las métricas obtenidas para los modelos Logistic Regression, SVM y XGBoost, tanto de sus versiones individuales como aquellas basadas en \textit{pipelines}. En esta se observa que los resultados obtenidos en las versiones correspondientes a \textit{pipelines} son, en general, superiores a las versiones individuales. En función de los valores obtenidos de ROC-AUC, f2-Score y \textit{recall} se seleccionan a las versiones basadas en \textit{pipelines} para cada uno de los modelos.

\begin{table}[H]
    \centering
    \caption{Resultados de métricas de modelos.}
    \begin{tabularx}{\textwidth}{c X c X X c X}
		\toprule
		\textbf{Modelo} & \textbf{Recall} & \textbf{Precision} & \textbf{F1-Score} & \textbf{F2-Score} & \textbf{AUPRC} & \textbf{ROC-AUC}\\
		\midrule
        Logistic Regression & 0,6818 & 0,1480 & 0,2432 & 0,3961 & 0,2321 & 0,9009 \\
        LR - Pipeline & 0,8636 & 0,1347 & 0,2331 & 0,4148 & 0,2615 & 0,9181 \\
        SVM & 0,0303 & 0,1428 & 0,05 & 0,0359 & 0,2158 & 0,8854 \\
        SVM - Pipeline & 0,8181 & 0,1436 & 0,2443 & 0,4218 & 0,2405 & 0,9101 \\
        XGBoost & 0,4242 & 0,3544 & 0,3862 & 0,4081 & 0,3327 & 0,9047 \\
        XGBoost - Pipeline & 0,7727 & 0,2170 & 0,3388 & 0,5110 & 0,3388 & 0,9236 \\
        \bottomrule
		\hline
    \end{tabularx}
    \label{tab:resultados-metricas}
\end{table}

\subsection{Comparación entre modelos}

Una vez obtenidas las métricas de los modelos, se procede a compararlos entre sí. Para ello, se hace uso de la métrica AUPRC o \textit{average precision}. Por este motivo, si se observa la tabla \ref{tab:resultados-metricas}, el mejor modelo es XGBoost - pipeline.

El \textit{model registry} de MLFlow se utiliza para seleccionar el modelo correcto. En este, se registran las distintas versiones del modelo final del trabajo, cuyo nombre es \texttt{bankruptcy\_prediction}. En sus metadatos se indica qué \textit{run} o ejecución lo generó. Tal como se observa en la figura \ref{fig:best-model}, en la versión número 4 del modelo final se detalla que se genera a partir de XGBoost, lo que indica que la comparación entre modelos funciona correctamente.

\begin{figure}[H]
	\centering
	\includegraphics[width=0.88\textwidth]{./Figures/best_model.png}
	\caption{Mejor modelo - \textit{Model registry} de MLFLow.}
	\label{fig:best-model}
\end{figure}

\subsection{Matriz de confusión normalizada del mejor modelo}

La figura \ref{fig:matriz-confusion-mejor-modelo} corresponde a la matriz de confusión normalizada por fila del mejor modelo (XGBoost - pipeline). Esta permite conocer qué tan bien clasifica el modelo a los casos de quiebra como así también a los casos de empresas sanas.

La matriz indica que el modelo obtiene un \textit{recall} o \textit{true positive rate} de $0,79$, es decir, que el modelo detecta de manera satisfactoria al $79\%$ de las empresas que entran en quiebra. El $13\%$ que resta se clasifica de manera errónea, es decir, el modelo indica que no entra en quiebra. Esto se ve en el valor de \textit{false negative rate}, que es de $0,21$.

En lo que respecta a las empresas sanas, el modelo detecta con éxito al $91\%$, tal como lo indica el \textit{true negative rate} o \textit{specifity} de $0,91$. El $9,4\%$ de casos que restan se clasifican de manera equivocada como quiebras, lo que se refleja en el \textit{false positive rate} de $0,094$.

En resumen, esta matriz describe de manera simple qué tan bien identifica el modelo a cada clase. Para el caso particular del dataset de Taiwán, la matriz indica que el modelo detecta el $79\%$ los casos de quiebra. Esta es una buena métrica, ya que el objetivo de este modelo es minimizar las quiebras no detectadas ($21\%$), debido al elevado riesgo financiero que supone un impago. Como contraparte, el modelo tiene una tasa de falsos positivos de $9,4\%$, que representan a empresas sanas marcadas como quiebra.

\begin{figure}[H]
	\centering
	\includegraphics[width=0.95\textwidth]{./Figures/fig_matriz_confusion.png}
	\caption{Matriz de confusión del mejor modelo.}
	\label{fig:matriz-confusion-mejor-modelo}
\end{figure}

En lo que respecta al valor de falsos negativos ($21\%$), el modelo podría clasificarlos de manera errónea por razones ajenas a las variables financieras, por ejemplo, cambio en regulaciones financieras del mercado de Taiwán, crisis inesperadas, entre otras.

\subsection{Curva ROC del mejor modelo}

La figura \ref{fig:fig_roc_curve} muestra la curva ROC del mejor modelo. Esta grafica el \textit{recall} o \textit{true positive rate} del modelo contra el \textit{false positive rate} en distintos umbrales o \textit{thresholds}, que muestran la disyuntiva de identificar de manera correcta un caso positivo y marcar de manera errónea los negativos. El área bajo esta curva (AUC-ROC) permite determinar qué tan bien el modelo separa las clases. Un valor cercano a $1$ indica que se tiene un buen modelo, mientras que un valor cercano a $0,5$ muestra que el modelo tiene un comportamiento azaroso.

\begin{figure}[H]
	\centering
	\includegraphics[width=0.92\textwidth]{./Figures/fig_roc_curve.png}
	\caption{Curva ROC del mejor modelo.}
	\label{fig:fig_roc_curve}
\end{figure}

Para el mejor modelo se tiene un AUC-ROC de $0,9236$, lo que indica que el modelo es capaz de distinguir empresas en quiebra de empresas sanas en diferentes umbrales.

Igualmente, esta curva y su métrica asociada no son suficientes para evaluar la performance del modelo. Esto se debe a que la curva hace uso de la métrica de \textit{false positive rate}, cuya fórmula se observa en la ecuación \ref{eq:false_positive_rate}, donde FP es \textit{false positive} y TN \textit{true negative}.

\begin{equation}
	\label{eq:false_positive_rate}
	FPR = \frac{FP}{FP + TN}
\end{equation}

El dataset de Taiwán tiene 6599 casos de empresas sanas o \textit{true negative}. Por lo tanto, incluso si el modelo genera cientos de falsos positivos, el valor de \textit{false positive rate} será cercano a cero, lo que hace que la curva ROC parezca excesivamente optimista.

\subsection{Curva \textit{precision-recall} del mejor modelo}

La curva \textit{precision-recall} \citep{WEBSITE:precision-recall-curve} es un método ideal para evaluar modelos que trabajan con datasets desbalanceados, en los que detectar la clase positiva minoritaria es crucial. Esto se debe a que, a diferencia de la curva ROC, no hace uso de la métrica de \textit{false positive rate}. Esta curva grafica distintos umbrales que muestran la disyuntiva de identificar correctamente las empresas que quiebran (\textit{recall}) y garantizar que estas predicciones sean precisas (\textit{precision}).

A diferencia de la curva ROC, el \textit{baseline} o comportamiento azaroso se define como la proporción de la clase positiva en el dataset. Para el caso del dataset de Taiwán, el baseline es de $0,32$ o $3,2\%$ ($220$ quiebras entre $6819$ empresas totales). El valor del área bajo la curva \textit{precision-recall} (AUPRC) del mejor modelo es $0,3388$ o $33,88\%$, 10 veces mejor que el \textit{baseline}. A partir de estos valores se concluye que el modelo extra patrones financieros de alto valor predictivo. Igualmente, la diferencia entre las métricas AUC-ROC y AUPRC muestra que, si bien el modelo clasifica la quiebra con éxito, sus predicciones positivas contienen un elevado volumen de falsas alarmas.

La figura \ref{fig:fig_precision_recall_curve} muestra la curva \textit{precision-recall} y el \textit{baseline} del mejor modelo (XGBoost - pipeline).

\begin{figure}[H]
	\centering
	\includegraphics[width=0.80\textwidth]{./Figures/fig_precision_recall_curve.png}
	\caption{Curva \textit{precision-recall} del mejor modelo.}
	\label{fig:fig_precision_recall_curve}
\end{figure}

\subsection{Variables más importantes del mejor modelo}

El modelo XGBoost, que se implementa con la clase \texttt{XGBClassifier} \citep{WEBSITE:xgboost_classifier} del paquete \texttt{XGBoost} \citep{WEBSITE:package_xgboost}, tiene el atributo \texttt{feature\_importance}, que indica qué tan útil fue cada \textit{feature} a la hora de configurar el modelo. Este utiliza por defecto la métrica de \textit{gain} o ganancia \citep{model:xgboost} y mide la contribución de una \textit{feature} según cuánto reduce el error total del modelo al dividir los datos en los árboles.

La tabla \ref{tab:var_importance_best_model} muestra las diez \textit{features} más importantes del modelo.

% \begin{table}[h]
% 	\centering
% 	\caption[Variables más importantes del mejor modelo]{Variables más importantes del mejor modelo.}
% 	\begin{tabular}{c c}
% 		\toprule
% 		\textbf{Feature} 	 & \textbf{Ganancia (\textit{gain})} \\
% 		\midrule
%         Working Capital/Equity &   0,163086 \\
%         Net Value Per Share (C) &   0,096922\\
%         Inventory Turnover Rate (times) &   0,069746 \\
%         Cash Flow to Total Assets &   0,062084 \\
%         Interest-bearing debt interest rate &   0,043152 \\
%         Tax rate (A) &   0,040591 \\
%         Quick Ratio &   0,029364 \\
%         Current Liability to Assets &   0,024506 \\
%         Net Value Per Share (B) &   0,021983 \\
%         Total income/Total expense &   0,017151 \\
% 		\bottomrule
% 		\hline
% 	\end{tabular}
% 	\label{tab:var_importance_best_model}
% \end{table}

% \begin{table}[h]
% 	\centering
% 	\caption[Variables más importantes del mejor modelo]{Variables más importantes del mejor modelo.}
% 	% \begin{tabular}{c c c}
%     \begin{tabularx}{\textwidth}{l X c}
% 		\toprule
% 		\textbf{Feature} & \textbf{Descripción}	 & \textbf{Ganancia (\textit{gain})} \\
% 		\midrule
%         Working Capital/Equity & Proporción del capital de trabajo financiado por los accionistas. &   0,163086 \\
%         Net Value Per Share (C) & Valor contable por acción bajo criterios estrictos de liquidación o consolidación reciente. &  0,096922\\
%         Inventory Turnover Rate (times) & Frecuencia anual con la que se vende y reemplaza el inventario. &  0,069746 \\
%         Cash Flow to Total Assets & Capacidad de los activos para generar efectivo neto. &  0,062084 \\
%         Interest-bearing debt interest rate & Costo financiero o tasa cobrada por los acreedores. &  0,043152 \\
%         Tax rate (A) &  Tasa impositiva efectiva que afecta los beneficios netos. & 0,040591 \\
%         Quick Ratio & Capacidad de pago inmediato sin liquidar inventarios corrientes. &  0,029364 \\
%         Current Liability to Assets &  Proporción de activos financiada con deuda a corto plazo. & 0,024506 \\
%         Net Value Per Share (B) & Valor en libros tradicional por acción circulante bajo principios contables estándar. &  0,021983 \\
%         Total income/Total expense & Índice de eficiencia operativa global y rentabilidad básica. &  0,017151 \\
% 		\bottomrule
% 		\hline
% 	% \end{tabular}
%     \end{tabularx}
% 	\label{tab:var_importance_best_model}
% \end{table}

\begin{table}[H]
	\centering
	\caption[Variables más importantes del mejor modelo]{Variables más importantes del mejor modelo.}
	% \begin{tabular}{c c c}
    \begin{tabularx}{\textwidth}{c c}
		\toprule
		\textbf{Feature} & \textbf{Ganancia (\textit{gain}) normalizada} \\
		\midrule
        \texttt{Working Capital/Equity} &   0,163086 \\
        \texttt{Net Value Per Share (C)} &  0,096922\\
        \texttt{Inventory Turnover Rate (times)} &  0,069746 \\
        \texttt{Cash Flow to Total Assets} &  0,062084 \\
        \texttt{Interest-bearing debt interest rate} &  0,043152 \\
        \texttt{Tax rate (A)} & 0,040591 \\
        \texttt{Quick Ratio} &  0,029364 \\
        \texttt{Current Liability to Assets} & 0,024506 \\
        \texttt{Net Value Per Share (B)} &  0,021983 \\
        \texttt{Total income/Total expense} &  0,017151 \\
		\bottomrule
		\hline
	% \end{tabular}
    \end{tabularx}
	\label{tab:var_importance_best_model}
\end{table}

% En la tabla \ref{tab:var_importance_best_model} se observa que la variable que mayor ganancia acumula es \texttt{Working Capital/Equity} con un $16,3\%$. Es decir, el modelo considera que una de las principales causas de quiebra de una empresa es una mala liquidez operativa, porque si el capital de trabajo que respaldan los accionistas se agota, la firma pierde su escudo operativo básico.

% La siguiente variable con mayor ganancia es \texttt{Net Value Per Share (C)}, con un valor de $9,69\%$. Esto indica que el modelo evalúa las caídas drásticas en el valor real de liquidación por acción para predecir una quiebra de una empresa.

En la tabla \ref{tab:var_importance_best_model} se observa que las variables \texttt{Working Capital/Equity} (representa la proporción del capital de trabajo financiado por los accionistas) y \texttt{Net Value Per Share (C)} (valor contable por acción bajo criterios estrictos de liquidación o consolidación reciente) acumulan cerca del $26\%$ de la ganancia, por lo que son las que más le permiten al modelo definir una quiebra. Esto podría indicar que el modelo asocia la quiebra con empresas que agotan el capital de trabajo operativo que respaldan sus socios y cuyos valores de liquidación por acción disminuyen o incluso se desploman.

Por su parte, las variables \texttt{Inventory Turnover Rate (times)} (consiste en la frecuencia anual con la que se vende y reemplaza el inventario), \texttt{Cash Flow to Total Assets} (representa la capacidad de los activos para generar efectivo neto) e \texttt{Interest-bearing debt interest rate} (costo financiero o tasa cobrada por los acreedores) acumulan cerca del $22\%$ de la ganancia del modelo. Esto podría indicar que el modelo penaliza a las empresas que no rotan sus inventarios y que no generan caja real con sus activos. Esta penalización sería aún mayor si las empresas deben pagar tasas altas de interés a sus acreedores.

En general, las variables más importantes del modelo podrían indicar que este determina la quiebra a partir de la liquidez inmediata y la eficiencia del ciclo de efectivo de una empresa dada.

%----------------------------------------------------------------------------------------
%	SECTION 2
%----------------------------------------------------------------------------------------
\section{Evaluación de entorno Airflow}
\label{sec:evaluacion_airflow}

En la sección \ref{sec:automatizacion_airflow} se explica cómo se automatizan procesos mediante Airflow. En esta sección se explica cómo se evalúa su correcto funcionamiento.

\subsection{DAG \texttt{download and split dataset}}

El primer DAG \citep{WEBSITE:airflow_dags} que se ejecuta es \texttt{download and split dataset}. Este se compone de dos tareas: \texttt{download\_dataset\_to\_s3}, que se encarga de descargar el dataset original y guardarlo en S3 (más precisamente en minio), y la tarea \texttt{split\_raw\_dataset\_and\_store\_to\_s3}, que se encarga de hacer la división en \textit{train} y \textit{test}, y guardar los resultados en S3/minio.

Entonces, para validar que se ejecuta correctamente se deben revisar los \textit{logs} de Airflow y revisar que los datasets estén siendo guardados en S3/minio.

En la figura \ref{fig:validacion_dag01} se observa la correcta ejecución del dag \texttt{download and split dataset}, en donde ambas tareas aparecen marcadas como correctas.

\begin{figure}[htbp]
	\centering
	\includegraphics[width=0.95\textwidth]{./Figures/download_and_split_dataset.png}
	\caption{Validación DAG \texttt{download and split dataset}.}
	\label{fig:validacion_dag01}
\end{figure}

En la figura \ref{fig:validacion_datasets_minio} se observan en S3/minio los datasets de \textit{raw}, que es generado por la tarea \texttt{download\_dataset\_to\_s3}, como así también los datasets de \textit{train} y \textit{test}, ambos generados por la tarea \texttt{split\_raw\_dataset\_and\_store\_to\_s3}.

\begin{figure}[H]
	\centering
	\includegraphics[width=0.80\textwidth]{./Figures/datasets_minio.png}
	\caption{Validación datasets en S3/minio.}
	\label{fig:validacion_datasets_minio}
\end{figure}

\subsection{DAG \texttt{model training}}

En la sección \ref{sec:model_training} se detalla al DAG de \texttt{model training}. Este está compuesto por tres tareas que se ejecutan en paralelo, \texttt{logistic\_regression\_training}, que entrena un modelo Logistic Regression, \texttt{svm\_training}, que entrena un modelo SVM, y \texttt{xgboost\_training}, que entrena un modelo XGBoost. Una vez finalizadas estas tareas, se ejecuta una tarea final, \texttt{seleccionar\_champion}, cuyo objetivo es el de seleccionar al mejor modelo.

Para validar que este DAG funciona correctamente se hace uso de los logs de Airflow. Además, se valida que cada tarea registre correctamente los modelos entrenados en MLFlow.

En la figura \ref{fig:validacion_dag02} se observan los resultados del DAG \texttt{model training}, en que se detalla que cada tarea tuvo un resultado satisfactorio.

\begin{figure}[htbp]
	\centering
	\includegraphics[width=0.87\textwidth]{./Figures/model_training_validation.png}
	\caption{Validación DAG \texttt{model training}.}
	\label{fig:validacion_dag02}
\end{figure}

%----------------------------------------------------------------------------------------
%	SECTION 3
%----------------------------------------------------------------------------------------
\section{Evaluación de selección del mejor modelo con otras métricas}

En la sección \ref{sec:pruebas_indiv_modelos} se detalla cómo se evalúa la selección del mejor modelo mediante el uso del \textit{model registry} de MLFlow. En este caso, se hace uso de la métrica AUPRC o \textit{average precision}, ya que es la mejor para casos en que el dataset presenta desbalance. Igualmente, el sistema desarrollado permite elegir otras métricas con las que se elige al mejor modelo.

Si se utiliza la métrica de \textit{recall}, los modelos quedan ordenados tal como se detalla en la tabla \ref{tab:comparacion-modelos-recall}. De esta manera, el modelo a elegir como \textit{champion} es "Logistic Regression - Pipeline".

\begin{table}[h]
	\centering
	\caption[Comparación mediante recall]{Comparación entre modelos mediante \textit{recall}.}
	\begin{tabular}{c c}
		\toprule
		\textbf{Modelo} 	 & \textbf{recall} \\
		\midrule
        Logistic Regression - Pipeline & 0,8636 \\
        SVM - pipeline & 0,8181 \\
        XGBoost - pipeline & 0,7727 \\
		\bottomrule
		\hline
	\end{tabular}
	\label{tab:comparacion-modelos-recall}
\end{table}

En la figura \ref{fig:best_model_recall} se observa que el nuevo mejor modelo, cuya versión es la número 5, fue creado en base al modelo Logistic Regression, tal como debe suceder al elegir la métrica \textit{recall} como método de decisión.

\begin{figure}[htbp]
	\centering
	\includegraphics[width=0.95\textwidth]{./Figures/best_model_recall.png}
	\caption{Best model - métrica recall.}
	\label{fig:best_model_recall}
\end{figure}

%----------------------------------------------------------------------------------------
%	SECTION 4
%----------------------------------------------------------------------------------------
\section{Evaluación de servicio web}
\label{sec:eval_servicio_web}

En esta sección se detallan las pruebas realizadas a cada uno de los \textit{endpoints} del servicio web. Estas pruebas se realizan mediante la herramienta Postman \citep{WEBSITE:postman}.

\subsection{Evaluación de \textit{endpoint} \texttt{get\_champion\_image}}
\label{sec:eval_get_champion_image}
El \textit{endpoint} \texttt{get\_champion\_image} permite obtener una imagen asociada al modelo \textit{champion}. Este se define mediante la URL \texttt{champion/image/<name>}, en donde \texttt{<name>} puede tomar alguno de los siguientes valores:

\begin{itemize}
    \item \texttt{confusion\_matrix}: devuelve el gráfico de la matriz de confusión del modelo \textit{champion}.
    \item \texttt{confusion\_matrix\_normalized\_by\_row}: devuelve el gráfico de la matriz de confusión normalizada por fila del modelo \textit{champion}.
    \item \texttt{precision\_recall\_curve}: devuelve el gráfico de la curva \textit{precision-recall} del modelo \textit{champion}.
    \item \texttt{roc\_curve}: devuelve el gráfico de la curva ROC del modelo \textit{champion}.
\end{itemize}

A partir de esto, se definen dos tipos de pruebas: una en la que se busca validar que, al ingresar uno de los nombres de imágenes correcto, se obtiene el gráfico correspondiente; y otra en la que, al ingresar un nombre erróneo, se obtiene un mensaje de error.

En la figura \ref{fig:champion_image_01} se observa el resultado obtenido al realizar una solicitud del gráfico de \texttt{confusion matrix normalized by row}. Esto se logra al utilizar el valor de \texttt{confusion\_matrix\_normalized\_by\_row} para el campo \texttt{<name>} de la URL de este \textit{endpoint}.

\begin{figure}[H]
	\centering
	\includegraphics[width=0.9\textwidth]{./Figures/champion_image_01.png}
	\caption{Best model - \texttt{Confusion matrix normalized by row}.}
	\label{fig:champion_image_01}
\end{figure}

En la figura \ref{fig:champion_image_02} se observa el resultado obtenido al utilizar un valor erróneo. En este caso, el valor que se utiliza es \textit{invalid}.

\begin{figure}[H]
	\centering
	\includegraphics[width=0.95\textwidth]{./Figures/champion_image_02.png}
	\caption{Best model - Imagen inválida.}
	\label{fig:champion_image_02}
\end{figure}

\subsection{Evaluación de \textit{endpoint} \texttt{get\_champion\_metrics}}
\label{sec:eval_get_champion_metrics}
El \textit{endpoint} \texttt{get\_champion\_metrics} permite obtener las métricas asociadas al mejor modelo o modelo \textit{champion}. Se define con la URL \texttt{champion/metrics}, y se caracteriza por no tener parámetros. En la figura \ref{fig:champion_metrics} se observan los resultados de evaluar este \textit{endpoint}.

\begin{figure}[htbp]
	\centering
	\includegraphics[width=0.90\textwidth]{./Figures/champion_metrics.png}
	\caption{Champion - metrics.}
	\label{fig:champion_metrics}
\end{figure}

\subsection{Evaluación de \textit{endpoint} \texttt{get\_experiment\_image}}

El \textit{endpoint} \texttt{get\_experiment\_image} se define a partir de la siguiente URL:

\begin{itemize}
    \item \texttt{experiment/<experiment>/image/<image>}
\end{itemize}

Esta tiene dos parámetros: \texttt{<experiment>}, que indica a qué experimento se quiere consultar, siendo \texttt{logistic\_regression}, \texttt{svm} y \texttt{xgboost} los valores posibles; y el parámetro \texttt{<image>}, que indica qué imagen se desea obtener, siendo los valores posibles los mismos que para el \textit{endpoint} \texttt{get\_champion\_image}, detallados en la sección \ref{sec:eval_get_champion_image}.

Los métodos para evaluar este endpoint son similares a los descriptos en la sección \ref{sec:eval_get_champion_image}, es decir, utilizando tanto valores válidos como inválidos para los parámetros \texttt{<experiment>} e \texttt{<image>}.

\subsection{Evaluación de \textit{endpoint} \texttt{get\_experiment\_metrics}}

El \textit{endpoint} \texttt{get\_experiment\_metrics} permite obtener métricas de un experimento dado. Está definido por la siguiente URL:

\begin{itemize}
    \item \texttt{experiment/<experiment>/metrics}
\end{itemize}

Esta tiene un solo parámetro: \texttt{<experiment>}. Este acepta únicamente los valores \texttt{logistic\_regression}, \texttt{svm} y \texttt{xgboost}.

La manera de evaluar este \textit{endpoint} es similar a la de \texttt{get\_champion\_metrics}, detallada en la sección \ref{sec:eval_get_champion_metrics}. La única diferencia es que ahora se debe indicar a qué experimento se desea consultar. Si el experimento indicado es correcto, se verá un resultado similar al de la sección \ref{sec:eval_get_champion_metrics}. En caso contrario, se muestra un mensaje de error, en que se indica que el parámetro ingresado es erróneo, y se listan los valores posibles a utilizar.

\subsection{Evaluación de \textit{endpoint} \texttt{predict}}

El \textit{endpoint} \texttt{predict} permite realizar predicciones con el modelo \textit{champion} a partir de una serie de datos ingresados por el usuario. Este \textit{endpoint} se define a partir de la URL \texttt{predict}. No tiene parámetros en la URL, pero sí acepta un \texttt{body} en su \textit{request}, que debe tener la información de una empresa dada, detallando un valor para las 95 variables del dataset.

A partir de esta información, se devuelve un mensaje que indica si la empresa entra en quiebra o no. En la figura \ref{fig:predict} se observan los resultados de una prueba en la que la empresa consultada no entra en quiebra.

\begin{figure}[H]
	\centering
	\includegraphics[width=0.87\textwidth]{./Figures/predict.png}
	\caption{Endpoint \texttt{predict}.}
	\label{fig:predict}
\end{figure}

%----------------------------------------------------------------------------------------
%	SECTION 5
%----------------------------------------------------------------------------------------
\section{Evaluación del entorno MLFlow}

La evaluación del entorno MLFlow, tanto su \textit{model registry} como su \textit{tracking server}, se hizo en función de las integraciones con Airflow y con el servicio web.

Airflow hace uso del \textit{tracking server} a la hora de registrar los entrenamientos de los modelos Logistic Regression, SVM y XGBoost, mientras que al \textit{model registry} lo utiliza para guardar el mejor modelo. Esto se refleja en el DAG \texttt{model\_training}, más precisamente, en \texttt{logistic\_regression\_training}, \texttt{svm\_training} y \texttt{xgboost\_training} para el \textit{tracking server}; y en \texttt{seleccionar\_champion} para el caso del \textit{model registry}. Como las pruebas de la sección \ref{sec:evaluacion_airflow} son satisfactorias, se puede decir que la integración de MLFlow y Airflow también lo es.

Respecto al servicio web, este hace uso del \textit{model registry} para obtener el modelo \textit{champion} en los \textit{endpoints} de \texttt{get\_champion\_image}, \texttt{get\_champion\_metrics} y \texttt{predict}. Por su parte, el \textit{tracking server} se utiliza para obtener modelos específicos en los \textit{endpoints} de \textit{get\_experiment\_image} y \textit{get\_experiment\_metrics}. Como las pruebas de la sección \ref{sec:eval_servicio_web} son satisfactorias, la integración entre MLFlow y el servicio web también lo es.

De esta manera se comprueba que el entorno MLFlow, tanto a nivel de \textit{tracking server} como de \textit{model registry}, funciona correctamente.

%----------------------------------------------------------------------------------------
%	SECTION 6
%----------------------------------------------------------------------------------------
% \section{Evaluación con otros dataset}

% A fin de determinar si el mejor modelo permite predecir quiebra de empresas de cualquier país, se opta por evaluarlo utilizando otros datasets.

% \subsection{Dataset polaco}
% \label{subsec:dataset_polaco}

% Se utiliza el dataset de \textit{Polish Companies Bankruptcy} \citep{dataset:polish_companies_bankruptcy_365}, que contiene información financiera de empresas polacas entre los años 2000 y 2012. La particularidad de este dataset es que tiene menos columnas que el dataset de Taiwán, ya que solo 11 son equivalentes. Estas son:
% % Estas se detallan en la tabla \ref{tab:polish-taiwan}.

% \begin{itemize}
%     \item \texttt{Net Income to Total Assets}: su equivalente en el dataset de Polonia es \texttt{Attr1 (net profit / total assets)}.
%     \item \texttt{Total debt/Total net worth}:  su equivalente en el dataset de Polonia es \texttt{Attr2 (total liabilities / total assets)}.
%     \item \texttt{Working Capital to Total Assets}: su equivalente en el dataset de Polonia es \texttt{Attr3 (working capital / total assets)}.
%     \item \texttt{Current Ratio}: la columna equivalente del dataset de Polonia es \texttt{Attr4 (current assets / short-term liabilities)}.
%     \item \texttt{Retained Earnings to Total Assets}: su equivalente en el dataset de Polonia es \texttt{Attr6 (retained earnings / total assets)}.
%     \item \texttt{ROA(C) before interest...}: su equivalente en el dataset de Polonia es \texttt{Attr7 (EBIT / total assets)}.
%     \item \texttt{Revenue Per Share (Yuan ¥}: su equivalente en el dataset de Polonia es \texttt{Attr9 (sales / total assets)}.
%     \item \texttt{ROA(B) before interest and depreciation after tax}: igual a \texttt{Attr11 (gross profit / total assets)} del dataset polaco.
%     \item \texttt{Pre-tax net Interest Rate}: su equivalente en el dataset de Polonia es \texttt{Attr14 (profit on operating activities / ...)}.
%     \item \texttt{Quick Ratio}: la columna equivalente del dataset de Polonia es \texttt{Attr46 (current assets - inventory / ...)}.
%     \item \texttt{Current Liability to Assets}: su equivalente en el dataset de Polonia es \texttt{Attr51 (short-term liabilities / total assets)}.
%     \item \texttt{Total income/Total expense}: su equivalente en el dataset de Polonia es \texttt{Attr58 (total costs / total sales)}.
% \end{itemize}

% % \begin{table}[H]
% %     \centering
% %     \caption{Equivalencia de columnas entre datasets de Polonia y Taiwán.}
% %     \begin{tabularx}{\textwidth}{X X}
% % 		\toprule
% % 		\textbf{Columna - Polonia} & \textbf{Columna - Taiwán} \\
% % 		\midrule
% %         \texttt{Attr1 (net profit / total assets)}			& \texttt{Net Income to Total Assets} \\
% %         \texttt{Attr2 (total liabilities / total assets)}	& \texttt{Total debt/Total net worth} \\
% %         \texttt{Attr3 (working capital / total assets)}		& \texttt{Working Capital to Total Assets} \\
% %         \texttt{Attr4 (current assets / short-term liabilities)}	& \texttt{Current Ratio} \\
% %         \texttt{Attr6 (retained earnings / total assets)}	& \texttt{Retained Earnings to Total Assets} \\
% %         \texttt{Attr7 (EBIT / total assets)}	& \texttt{ROA(C) before interest...} \\
% %         \texttt{Attr9 (sales / total assets)}	& \texttt{Revenue Per Share (Yuan ¥)} \\
% %         \texttt{Attr11 (gross profit / total assets)}	& \texttt{ROA(B) before interest and depreciation after tax} \\
% %         \texttt{Attr14 (profit on operating activities / ...)}	& \texttt{Pre-tax net Interest Rate} \\
% %         \texttt{Attr46 (current assets - inventory / ...)}	& \texttt{Quick Ratio} \\
% %         \texttt{Attr51 (short-term liabilities / total assets)}	& \texttt{Current Liability to Assets} \\
% %         \texttt{Attr58 (total costs / total sales)}	& \texttt{Total income/Total expense} \\
% % 		\bottomrule
% % 		\hline
% %     \end{tabularx}
% %     \label{tab:polish-taiwan}
% % \end{table}

% En cuanto a registros, el dataset de Polonia cuenta con 7027, de los que 6756 corresponden a empresas que no entran en quiebra, mientras que los 271 restantes corresponden a empresas que sí entran en quiebra.

% Para evaluar el dataset de Polonia, se realiza una conversión de nombres para que coincidan con los del dataset de Taiwán. En cuanto a las columnas faltantes, estas se dejan con valores iguales a cero, a fin de que el modelo se encargue de procesarlas. La evaluación se realiza con el mejor modelo, XGBoost con \textit{pipelines}, dando como resultado los valores que se muestran en la tabla \ref{tab:evaluacion-dataset-polonia}.

% \begin{table}[H]
% 	\centering
% 	\caption[Evaluación del dataset de Polonia]{Evaluación del dataset de Polonia.}
% 	\begin{tabular}{c c c}
% 		\toprule
% 		\textbf{Métrica} 	 & \textbf{Dataset de Polonia} & \textbf{Dataset de Taiwán} \\
% 		\midrule
%         recall & 1 & 0,7727 \\
%         precision & 0,0386 & 0,2227 \\
%         f1-score & 0,0743 & 0,3457 \\
%         f2-score & 0,1672 & 0,5172 \\
%         AUPRC (average precision) & 0,0429 & 0,3259\\
%         AUC-ROC & 0,5096 & 0,9210 \\
% 		\bottomrule
% 		\hline
% 	\end{tabular}
% 	\label{tab:evaluacion-dataset-polonia}
% \end{table}

% Los valores correspondientes al dataset de Polonia son muy bajos. El valor de AUC-ROC, $0,5096$, indica que el comportamiento del modelo con estos datos es similar a un modelo aleatorio. El \textit{recall} es de $1$, es decir, clasifica correctamente a todos los casos de quiebra, pero al analizar la \textit{precision}, cuyo valor es $0,0386$, se deduce que el modelo clasifica la mayoría de los casos como quiebra, independientemente de si lo son o no. Esta diferencia de performance puede deberse a varios motivos, tales como:

% \begin{itemize}
%     \item La información presente en cada dataset: el dataset de Polonia tiene 11 columnas, mientras que el de Taiwán tiene 95.
%     \item A cómo se obtuvieron los valores.
%     \item A cómo se mide la quiebra en estos países: pueden tener diferentes normas contables y legislaciones de quiebra.
% \end{itemize}

% \subsection{Dataset ruso}
% \label{subsec:dataset_ruso}

% Se utiliza el dataset de RFSD o \textit{The Russian Financial Statements Database} \citep{dataset:institute_for_the_rule_of_law_at_european_university_at_st._petersburg_2025}, que contiene información financiera de empresas rusas entre los años 2011 y 2024. En particular, se utiliza un subconjunto de un millón de registros correspondientes al año 2024. No se utiliza el total de registros (tres millones) por limitaciones de la memoria RAM del entorno de trabajo.

% De manera similar al dataset polaco (ver sección \ref{subsec:dataset_polaco}), no todas las columnas coinciden con el dataset de Taiwán. Por ello, las faltantes son cargadas con valores iguales a cero, mientras que las columnas equivalentes se obtienen de la siguiente manera:

% \begin{itemize}
%     \item \texttt{ROA(C) before interest}: división entre las columnas \texttt{line\_2400} y \texttt{line\_1600} del dataset ruso. Representa la utilidad neta sobre el activo total.
%     \item \texttt{Current Ratio}: división entre las columnas \texttt{line\_1200} y \texttt{line\_1500} del dataset ruso. Representa la capacidad de pago a corto plazo.
%     \item \texttt{Quick Ratio}: representa la liquidez inmediata. Se calcula mediante la fórmula ${(\texttt{line\_1200} - \texttt{line\_1210}) / \texttt{line\_1500}}$.
%     \item \texttt{Net worth/Assets}: división de las columnas \texttt{line\_1300} y \texttt{line\_1600} del dataset ruso. Es la proporción de capital propio en la empresa.
%     \item \texttt{Debt ratio \%}: división entre las columnas \texttt{line\_1500} y \texttt{line\_1600} del dataset ruso. Consiste en el nivel de apalancamiento.
%     \item \texttt{Inventory Turnover Rate}: división entre las columnas \texttt{line\_2120} y \texttt{line\_1210} del dataset ruso. Consiste en la eficiencia en la rotación de mercancía.
%     \item \texttt{Current Assets/Total Assets}: división de las columnas \texttt{line\_1200} y \texttt{line\_1600} del dataset ruso. Representa el peso de los activos líquidos en el balance.
%     \item \texttt{Inventory/Total Assets}: división entre las columnas \texttt{line\_1210} y \texttt{line\_1600} del dataset ruso. Representación del peso de las existencias en el activo total.
%     \item \texttt{Operating Revenue}: columna \texttt{line\_2110} del dataset ruso. Son los ingresos brutos por ventas/operación.
%     \item \texttt{Net Income}: columna \texttt{line\_2400} del dataset ruso. Representa el resultado final del ejercicio.
% \end{itemize}

% Otra diferencia con el dataset de Taiwán es la columna de bancarrota. El dataset ruso no indica si una empresa quiebra o no, sino que esto se calcula con las siguientes columnas:

% \begin{itemize}
%     \item \texttt{dissolution\_date}: indica la fecha de disolución de una empresa. Si no es nula, la empresa sí entró en quiebra.
%     \item \texttt{filed}: indica si una empresa presenta estados financieros. Si toma el valor $0$, la empresa sí entró en quiebra.
% \end{itemize}

% La evaluación se realiza con el mejor modelo, XGBoost con \textit{pipelines}, dando como resultado los valores que se muestran en la tabla \ref{tab:evaluacion-dataset-ruso}.

% \begin{table}[H]
% 	\centering
% 	\caption[Evaluación del dataset ruso]{Evaluación del dataset ruso.}
% 	\begin{tabular}{c c c}
% 		\toprule
% 		\textbf{Métrica} 	 & \textbf{Dataset de Rusia} & \textbf{Dataset de Taiwán} \\
% 		\midrule
%         recall & 1 & 0,7727 \\
%         precision & 0,2993 & 0,2227 \\
%         f1-score & 0,4606 & 0,3457 \\
%         f2-score & 0,6810 & 0,5172 \\
%         AUPRC (average precision) & 0,6422 & 0,3259\\
%         AUC-ROC & 0,8773 & 0,9210 \\
% 		\bottomrule
% 		\hline
% 	\end{tabular}
% 	\label{tab:evaluacion-dataset-ruso}
% \end{table}

% En este caso se obtuvieron mejores métricas que con el dataset de Polonia. Algunas de estas son incluso mejores que la del dataset original. La métrica AURC-ROC, con valor de $0,8773$, es mucho mejor que el caso del dataset polaco ($0,5096$, que indica un comportamiento azaroso), pero menor a la del dataset de Taiwán ($0,9210$). El \textit{recall} es $1$, pero a diferencia del dataset polaco, la métrica de \textit{precision} es de $0,2993$. Esto indica que el modelo genera muchos falsos positivos. Por ejemplo, de $10$ casos clasificados como quiebra, solo $3$ serán correctos. Por su parte, la métrica AUPRC ($0,6422$) es mejor que las obtenidas en el dataset de Polonia ($0,0429$) y de Taiwán ($0,3259$). Esto indica que el modelo mantiene una buena relación entre \textit{precision} y \textit{recall}.

% Se puede decir que el modelo entrenado con el dataset de Taiwán puede detectar quiebras asociadas al mercado ruso, pero es sensible a clasificar erróneamente a empresas sanas. Esta imprecisión puede deberse al formato del dataset ruso, que solamente tiene $11$ columnas equivalentes con el dataset de Taiwán. De manera similar al caso Polonia, pueden haber también diferentes legislaciones y/o definiciones que inciden en la definición de una quiebra.

% \subsection{Dataset de Estados Unidos}
% \label{subsec:dataset_usa}

% Se utiliza el dataset de \textit{Bankruptcy prediction dataset for american companies in the stock market} \citep{dataset:usa_bankruptcy}, que contiene información financiera de empresas que cotizan en bolsas estadounidenses (NYSE y NASDAQ). Esta información cubre los años entre 1998 y 2018, y se divide en 78682 y 21 columnas. Solo 13 de estas 21 columnas tienen un equivalente en el dataset de Taiwán:

% \begin{itemize}
%     \item \texttt{Bankrupt?}: se obtiene a partir de la columna \texttt{status\_label} del dataset de Estados Unidos, y se trata de la variable objetivo. Se caracteriza por estar desbalanceada (5220 presentan quiebra, mientras que 73462 no).
%     \item \texttt{Net Income to Total Assets}: división entre las columnas \texttt{X6} y \texttt{X10} del dataset de Estados Unidos. Representa la rentabilidad sobre activos (ROA).
%     \item \texttt{Debt ratio \%}: división entre las columnas \texttt{X17} y \texttt{X10} del dataset de Estados Unidos. Representa la proporción de deuda.
%     \item \texttt{Current Ratio}: división entre las columnas \texttt{X1} y \texttt{X14} del dataset de Estados Unidos. Es la capacidad de la empresa para pagar sus deudas a corto plazo con sus activos corrientes.
%     \item \texttt{Retained Earnings to Total Assets}: división entre las columnas \texttt{X15} y \texttt{X10} del dataset de Estados Unidos. Mide el apalancamiento histórico y la capacidad de reinversión acumulada.
%     \item \texttt{Operating Gross Margin}: división entre las columnas \texttt{X13} y \texttt{X9} del dataset de Estados Unidos. Representa el margen bruto operativo.
%     \item \texttt{Net worth/Assets}: patrimonio neto sobre activos. Se calcula mediante la fórmula \texttt{(X10 - X17) / X10}.
%     \item \texttt{Quick Ratio}: se calcula mediante la fórmula \texttt{(X1 - X5) / X14}. Es la capacidad de pago inmediato.
%     \item \texttt{Liability to Equity}: consiste en la relación deuda/capital. Se calcula mediante la fórmula \texttt{X17 / (X10 - X17)}.
%     % \item \texttt{Net Income to Stockholder's Equity}: rentabilidad sobre capital. Se calcula mediante la fórmula \texttt{X6 / (X10 - X17)}.
%     \item \texttt{Net Income to Stockholder's Equity}: se calcula mediante la fórmula \texttt{X6 / (X10 - X17)}. Representa la rentabilidad sobre capital.
%     \item \texttt{Current Liability to Assets}: pasivo circulante sobre activos. Se obtiene al dividir las columnas \texttt{X14} y \texttt{X10} del dataset de Estados Unidos.
%     \item \texttt{ROA(C) before interest and ...}: rentabilidad operativa. Se obtiene al dividir las columnas \texttt{X4} y \texttt{X10}.
%     \item \texttt{Working Capital to Total Assets}: capital de trabajo sobre activos. Se calcula mediante la fórmula \texttt{(X1 - X14) / X10}.
% \end{itemize}

% La evaluación se realiza con el mejor modelo, XGBoost con \textit{pipelines}. Los resultados se muestran en la tabla \ref{tab:evaluacion-dataset-usa}.

% \begin{table}[H]
% 	\centering
% 	\caption[Evaluación del dataset estadounidense]{Evaluación del dataset estadounidense.}
% 	\begin{tabular}{c c c}
% 		\toprule
% 		\textbf{Métrica} 	 & \textbf{Dataset de Rusia} & \textbf{Dataset de Taiwán} \\
% 		\midrule
%         recall & 0,9941 & 0,7727 \\
%         precision & 0,0667 & 0,2227 \\
%         f1-score & 0,1250 & 0,3457 \\
%         f2-score & 0,2628 & 0,5172 \\
%         AUPRC (average precision) & 0,0663 & 0,3259\\
%         AUC-ROC & 0,4999 & 0,9210 \\
% 		\bottomrule
% 		\hline
% 	\end{tabular}
% 	\label{tab:evaluacion-dataset-usa}
% \end{table}

% Similar al caso del dataset polaco que se detalla en la sección \ref{subsec:dataset_polaco}, se obtuvieron métricas muy bajas. Se tiene un AUC-ROC de $0,4999$, que indica un comportamiento similar a un modelo aleatorio. Si bien el \textit{recall} es un valor muy alto ($0,9941$), si se analiza la \textit{precision} ($0,0667$), se concluye que el modelo clasifica a la mayoría de los casos como quiebra, independientemente de que lo sean o no.

% Esta baja performance puede deberse a:
%  \begin{itemize}
%      \item Diferencia de columnas entre datasets: el dataset de Taiwán tiene 95 columnas, pero el de Estados Unidos solo tiene 13 columnas equivalentes.
%      \item Diferentes formas de medir una quiebra: ambos países pueden tener diferentes normas contables y/o legislaciones de quiebra.
%  \end{itemize}